\subsection{Equipment Classification and Prediction Error}
The largest change in average absolute prediction error resulted from including the rental indicator alongside the HCPCS code. The actual-payment MAE decreased from \$47.84 to \$19.89 under the same 51-area holdout protocol. This suggests that averaging across rental statuses within a code discards predictive information relevant to the aggregate payments in this dataset. The improvement is attributable to a more specific descriptive baseline, not to a learned nonlinear representation. The code/rental grouping is also consistent with the unit of comparison used in the geographic coefficient-of-variation analysis.

\subsection{Standardized Targets Are Different Outcomes}
Standardized average payments had lower MAE and pooled NMAE than observed average payments for both baselines. Under rental-aware prediction, the two micro MAEs were \$10.30 and \$19.89. This direction was consistent across all 51 areas. The finding is compatible with the motivation for Medicare payment standardization described by O'Donnell et al.~\citep{odonnell2013}: some geographic payment variation is attributable to reimbursement adjustments rather than to the underlying service mix. Nevertheless, the two errors measure predictions of \emph{different target fields}. A smaller error for standardized payments does not establish that predicting the actual Medicare payment has improved. We did not evaluate a procedure that predicts standardized amounts and then transforms them back into observed payments.

\subsection{Uneven Geographic Reliability}
Pooled accuracy obscures geographic heterogeneity. For example, under the HCPCS-only actual-payment baseline, the macro MAE was \$46.27 while the highest area MAE was \$103.55 in Texas. The latter value cannot be interpreted as evidence of an intrinsic disadvantage associated with Texas: differences in HCPCS composition, rental mix, and availability of codes across training areas may contribute. The preliminary 30-seed comparison also did not show geographically grouped holdout consistently exceeding random-split error, emphasizing the need to examine both the evaluation design and test-set composition \citep{roberts2017}.

\subsection{Rare Codes and Residual Errors}
Unseen HCPCS codes comprised approximately 0.38\% of held-out observations under the HCPCS-only baseline but generated about 5.86\% of its total absolute error. The fallback to a global mean can be far from the appropriate amount for an unusual code. This is a useful diagnostic for model deployment because average accuracy may conceal occasional large errors. We have not established whether the worst errors arise from rare equipment types, claim aggregation, geographic payment adjustment, or other factors; such explanations require additional analysis.

\subsection{Limitations and Future Work}
The study uses one year of publicly released aggregate CMS records, which can reflect disclosure suppression and differences in service volume \citep{cmsdme2023}. The measured MAE is an unweighted average across aggregated rows; it does not represent beneficiary-level error or total expenditure error. The same payment-code hierarchy was used across targets, but target magnitudes and distributions are not identical, so percentage error reductions should not be interpreted as causal estimates of standardization's effect. In addition, no independent future-year test, statistical uncertainty interval for geographic comparisons, or flexible machine-learning model benchmark was completed. Dataset shift across time or settings remains a concern \citep{subbaswamy2020}. Future work could repeat the analysis across years, use service-volume-weighted metrics alongside unweighted metrics, evaluate within-code residual geography effects, and compare transparent baselines with more flexible estimators using equivalent geographic splits.